\documentclass[12pt,letterpaper]{article}
\usepackage{amsmath}
\usepackage{amsthm}
\usepackage{amsfonts}
\usepackage{amssymb}
\usepackage{amscd}
\usepackage{enumerate}
\usepackage{fancyhdr}
\usepackage{mathrsfs}
\usepackage{bbm}
\usepackage{framed}
\usepackage{mdframed}
\usepackage{cancel}
\usepackage{mathtools}
\usepackage{verbatim}
\usepackage{hyperref}
\usepackage[letterpaper,voffset=-.5in,bmargin=3cm,footskip=1cm]{geometry}
\setlength{\parindent}{0.0in}
\setlength{\parskip}{0.1in}
\allowdisplaybreaks
\headheight 15pt
\headsep 10pt
\newcommand\N{\mathbb N}
\newcommand\Z{\mathbb Z}
\newcommand\R{\mathbb R}
\newcommand\Q{\mathbb Q}
\newcommand\lcm{\operatorname{lcm}}
\newcommand\setbuilder[2]{\ensuremath{\left\{#1\;\middle|\;#2\right\}}}
\newcommand\E{\operatorname{E}}
\newcommand\V{\operatorname{V}}
\newcommand\Pow{\ensuremath{\operatorname{\mathcal{P}}}}

\DeclarePairedDelimiter\ceil{\lceil}{\rceil}
\DeclarePairedDelimiter\floor{\lfloor}{\rfloor}
\newcommand\hint[1]{\textbf{Hint}: #1}
\newcommand\note[1]{\textbf{Note}: #1}
\newcounter{enum22i}
\newenvironment{22enumerate}{\begin{enumerate}[a.]\itemsep0em\setcounter{enumi}{\value{enum22i}}}{\setcounter{enum22i}{\value{enumi}}\end{enumerate}}
\newenvironment{22itemize}{\begin{itemize}\itemsep0em}{\end{itemize}}
\fancypagestyle{firstpagestyle} {
  \renewcommand{\headrulewidth}{0pt}
  \lhead{\textbf{CSCI 0220}}
  \chead{\textbf{Discrete Structures and Probability}}
  \rhead{M. Littman}
}

\fancypagestyle{fancyplain} {
  \renewcommand{\headrulewidth}{0pt}
  \lhead{\textbf{CSCI 0220}}
  \chead{Final}
  \rhead{\textit{ }}
}
\pagestyle{fancyplain}
\usepackage{tikz}

\begin{document}
  \thispagestyle{firstpagestyle}
  \begin{center}
    {\huge \textbf{CSCI0220 Final}}

    {\large Due: \textit{Tuesday, August 10th at 11:59pm EDT}}
  \end{center}
  
 {\large \textbf{Please read in full:}}
 
You are not allowed to use the internet outside of the CSCI0220 website. Doing so is a breach of the \textbf{Collaboration Policy} and \textbf{Brown's Academic Code}.

You can refer to lecture slides, lecture captures, your own notes, the course textbook, the website resources, and Campuswire. You must complete the entire exam by yourself, without discussion of any kind with anyone else.

\textbf{By submitting this exam, you are confirming that you have understood these policies, and that you have followed them.}

As always, please do not include any identifying information about yourself in the handin, including your Banner ID. Be sure to fully explain your reasoning and show all work for full credit.

\subsection*{Problem 1}
{

{\setcounter{enum22i}{0}
	Suppose $R$ is a relation on a set $A$ and $q$ is a property of relations (such as
	reflexivity, symmetry, or transitivity). For property $q$, define
	the $q$-\textit{closure} of $R$ as the smallest relation on $A$ that has property $q$
	and is a superset of $R$. Recall that a set $S$ is a \textit{superset} of $T$ if and only if $T$ is
	a subset of $S$.

	For example, let $A = \{1,2,3\}$ and $R = \{(2,1)\}$. The \textit{reflexive closure} of R
	is the relation $\{(2,1),(1,1),(2,2),(3,3)\}$ while the \textit{symmetric closure} of R is
	the relation $\{(2,1),(1,2)\}$.

	\begin{22enumerate}

	\item Given a relation $R$, prove whether or not the transitive closure of the symmetric closure of the reflexive closure of $R$ is an equivalence relation. (That is, if one took the reflexive closure of $R$, and then the symmetric closure of the result, and then took the transitive closure of that, is the result an equivalence relation?)

	\item Provide a counterexample to show that, given a relation $R$, the reflexive closure of the symmetric closure of the transitive closure of $R$ is not always an equivalence relation. Justify your answer. (That is, find an $R$ such that if one took the transitive closure of $R$, and then the symmetric closure of the result, and then took the reflexive closure of that, the result is not an equivalence relation.)
	
	\hint{$|R|=2$ suffices to create a counterexample.}

	\end{22enumerate}

}
}

\newpage
\subsection*{Problem 2}
{
\nopagebreak 

{\setcounter{enum22i}{0}


	Let $X$ be a set of size $n$.  A partition of $X$ into $k$ blocks is defined as:
	$$X= \{B_1,B_2,\ldots,B_k\}$$
	such that
	\begin{center}
	\begin{enumerate}[1.]
	\item For all $i, B_i \neq \emptyset$.
	\item For all $i \neq j, B_i \cap B_j = \emptyset$.
	\item $B_1 \cup B_2 \cup \ldots \cup B_k = X$.
	\end{enumerate}
	\end{center}

	For all $0 < k \leq n$ define $S(n,k)$ to be the number of partitions of an $n$-element set into $k$ blocks.
	\begin{22enumerate}
	 \item Argue via counting argument that the following recurrence identity holds for integers $0 < k < n$ such that $n > 2, k \geq 2$:
	$$S(n,k) = kS(n-1,k) + S(n-1,k-1).$$
	\hint{Consider an arbitrary element. Where in the partition can this element be?}
	
	\item How many surjective functions from an $n$-element set to a $k$-element set are there in terms of $S(n,k)$, $n$, and $k$? Note that here $0 < k < n$. Justify your answer.
	
	\end{22enumerate}


}
}
\newpage
\subsection*{Problem 3}
{
{\setcounter{enum22i}{0}

    A \textit{bipartite graph} is a graph whose vertices can be partitioned into two disjoint vertex sets, $L$ and $R$, such that every edge has one endpoint in $L$ and the other in $R$. See the figure below for an example.
      \begin{figure}[h!]
        \caption{Three-vertex set $L$ (left) and two-vertex set $R$ (right)}
        \centering
          \includegraphics[scale = 0.3]{bipartite_graph.png}
      \end{figure}

      Take two integers $a,b \in \Z^+$ such that $a < b$, and create a bipartite graph where the vertex sets are $L=R=\{n \in Z \;|\; n \in [0, b-1]\}$. (We're imagining that the $0$ node in $L$ is distinct from the $0$ node in $R$, but they have the same name for easy reference.)

      Using edges, connect $i \in L$ to both $i \in R$ and $\mathrm{rem}(i+a,b) \in R$. An example is given below for $a=2$, $b=3$.
      
      \begin{figure}[h!]
        \caption{Example for $a=2$ and $b=3$.}
        \centering
          \includegraphics[scale = 0.4]{mod_graph.png}
      \end{figure}
      
      \begin{22enumerate}
          \item Under what conditions will there be a path from $0 \in L$ to $1 \in R$ and what will the minimal path length be, if it exists? Justify your answer.

          \item Consider now the graph where we connect vertex $i \in L$ to $\mathrm{rem}(i+ka,b) \in R$ for all $k \in \Z$. If we let $a=1$, find the expected minimal distance between two distinct, randomly selected vertices with respect to $b$, justifying your answer. No need to simplify your expression.
          
      \end{22enumerate}
     
    
}
}

\newpage

\subsection*{Problem 4}

{
{\setcounter{enum22i}{0}
Jupiter spotted this false proof of the `Graph Pigeonhole Principle.'\\

\textbf{Claim:} For $n\ge 4$, if an $n$-vertex graph has at least $n+1$ edges, it must contain a subgraph that is a triangle (a complete graph on $3$ vertices).

\textbf{Proof:}

Define $P(n)$: For every graph on $n$ vertices, if the graph has at
least $n+1$ edges, it contains a subgraph that is a triangle. 

Consider the base case $P(4)$. The complete graph $K_4$ has 6 edges and vertices $\{v_1, v_2, v_3, v_4\}$. We're interested in a graph on these vertices with at least $n+1=5$ edges. If the graph has 6 edges, it is a complete graph and any 3 of its vertices forms a triangle. If the graph has 5 edges, we can assume without loss of generality that the missing edge is
$(v_1,v_2)$. The edges $\{ v_1, v_3 \}$, $\{ v_3, v_4 \}$, and $\{
v_1, v_4 \}$ are in the graph and form a triangle. Hence, $P(4)$ is true, which proves our base case.

Assume that $P(k)$ is true for $n=k$. For the inductive step, consider a graph $G$ with $k$ vertices and at least $k+1$ edges. By our inductive hypothesis, $G$ contains a triangle. Construct a $(k+1)$-vertex graph $G'$ by adding one vertex to $G$ and any number of edges needed to ensure the number of edges is at least $k+2$. By the inductive hypothesis, $G$ contains a triangle. Since $G$ is a subgraph of $G'$, $G'$ also contains a triangle. Hence, if $P(k)$ is true, so is $P(k+1)$.

We have shown that $P(4)$ holds and that if $P(k)$ holds so does  $P(k+1)$. Thus, we have shown that $P(n)$ holds for all $n\ge 4$ by the principle of induction. QED.
\begin{22enumerate}
    \item Prove that Jupiter's Graph Pigeonhole Principle is false.
    
    \item What went wrong in Jupiter's proof?
    
    \item What is the expected value of the number of triangles formed in an $n$-vertex graph given that it has $n+1$ edges (selected uniformly at random)?
    
    \item Repair the theorem by showing that a graph with $n$ vertices where the minimum degree is more than $n/2$ must contain a triangle.
    
    \hint{Use the actual Pigeonhole Principle! Consider two connected vertices and their minimal degree.}
    
\end{22enumerate}
}
}


\end{document}